\chapter{Vereinigung und Schnitt}

\section{Elementkriterien und Inklusion}

\begin{definition}[Binäre Vereinigung und binärer Schnitt]
\[
 A\cup B=\bigcup\{A,B\},\qquad
 A\cap B=\{x\in A\mid x\in B\}.
\]
Es gelten die Elementkriterien
\[
 x\in A\cup B\Longleftrightarrow x\in A\lor x\in B,\qquad
 x\in A\cap B\Longleftrightarrow x\in A\land x\in B.
\]
\end{definition}

\begin{theorem}[Enthaltensein in Vereinigung und Schnitt]
\[
\begin{gathered}
 A\cap B\subseteq A,\quad A\cap B\subseteq B,\quad
 A\subseteq A\cup B,\quad B\subseteq A\cup B,\\
 C\subseteq A\cap B
       \Longleftrightarrow C\subseteq A\land C\subseteq B,\\
 A\cup B\subseteq C
       \Longleftrightarrow A\subseteq C\land B\subseteq C.
\end{gathered}
\]
\end{theorem}

\begin{theorem}[Inklusion durch Mengengleichungen]
\[
 A\subseteq B
   \Longleftrightarrow A\cap B=A
   \Longleftrightarrow A\cup B=B.
\]
\end{theorem}

\begin{theorem}[Vergrößern der beteiligten Mengen]
Für $A\subseteq C$ und $B\subseteq D$ gelten
\[
 A\cup B\subseteq C\cup D,\qquad A\cap B\subseteq C\cap D.
\]
\end{theorem}

\section{Grundlegende Mengengleichungen}

\begin{theorem}[Vertauschen, Klammern und Wiederholen]
\[
\begin{aligned}
 A\cup B&=B\cup A,& A\cap B&=B\cap A,\\
 (A\cup B)\cup C&=A\cup(B\cup C),&
 (A\cap B)\cap C&=A\cap(B\cap C),\\
 A\cup A&=A,& A\cap A&=A,\\
 A\cup\varnothing&=A,& A\cap\varnothing&=\varnothing.
\end{aligned}
\]
\cite{rutgers}
\end{theorem}

\begin{theorem}[Absorptionsgesetze]
\[
 A\cup(A\cap B)=A,\qquad A\cap(A\cup B)=A.
\]
\end{theorem}

\begin{theorem}[Distributivgesetze]
\[
\begin{aligned}
 A\cap(B\cup C)&=(A\cap B)\cup(A\cap C),\\
 A\cup(B\cap C)&=(A\cup B)\cap(A\cup C).
\end{aligned}
\]
\end{theorem}

\begin{theorem}[Gleichheit aus gemeinsamem Schnitt und gemeinsamer Vereinigung]
\[
 (A\cap C=B\cap C\ \land\ A\cup C=B\cup C)
       \Longleftrightarrow A=B.
\]
\end{theorem}

\section{Disjunktheit}

\begin{definition}[Disjunkte Mengen]
$A$ und $B$ heißen \emph{disjunkt}, wenn
$A\cap B=\varnothing$ gilt.
\end{definition}

\begin{theorem}[Elementkriterium und Vererbung der Disjunktheit]
\[
 A\cap B=\varnothing
     \Longleftrightarrow\neg\exists x\,(x\in A\land x\in B).
\]
Sind $A,B$ disjunkt sowie $C\subseteq A$ und $D\subseteq B$,
so sind auch $C,D$ disjunkt.
\end{theorem}

\begin{theorem}[Disjunktheit gegenüber einer Vereinigung]
\[
 A\cap(B\cup C)=\varnothing
 \Longleftrightarrow
 A\cap B=\varnothing\ \land\ A\cap C=\varnothing.
\]
Ferner gilt $A\cap A=\varnothing\Longleftrightarrow A=\varnothing$.
\end{theorem}
